\section{Espacios anillados}
\label{section:0.4-fr}

\subsection{Espacios anillados, $\mathcal{A}$-Módulos, $\mathcal{A}$-Álgebras}
\label{subsection:0.4.1-fr}

\begin{env}[4.1.1]
\label{0.4.1.1-fr}
Un \emph{espacio anillado} (resp. \emph{topológicamente anillado}) es un par
$(X,\mathcal{A})$ formado por un espacio topológico $X$ y un haz de anillos
(no necesariamente conmutativos) (resp. un haz de anillos topológicos)
$\mathcal{A}$ sobre $X$. Se dice que $X$ es el espacio
\oldpage[0\textsubscript{I}]{36}
topológico \emph{subyacente} al espacio anillado $(X,\mathcal{A})$, y que
$\mathcal{A}$ es el \emph{haz estructural}. Este último se denota por
$\mathcal{O}_X$, y su tallo en un punto $x\in X$ se denota por $\mathcal{O}_{X,x}$
o simplemente por $\mathcal{O}_x$ cuando no haya lugar a confusión.

Se designará por $1$ o $e$ la \emph{sección unidad} de $\mathcal{O}_X$
sobre $X$ (elemento unidad de $\Gamma(X,\mathcal{O}_X)$).

Como en este Tratado consideraremos sobre todo haces de anillos
\emph{conmutativos}, se sobreentenderá, cuando hablemos de un espacio
anillado $(X,\mathcal{A})$ sin más precisión, que $\mathcal{A}$ es un haz
de anillos conmutativos.

Los espacios anillados cuyo haz estructural no es necesariamente conmutativo
(resp. los espacios topológicamente anillados) forman una \emph{categoría}
si se define un \emph{morfismo}
$(X,\mathcal{A})\to(Y,\mathcal{B})$ como un par
$(\psi,\theta)=\Psi$ formado por una aplicación continua $\psi:X\to Y$ y un
$\psi$-\emph{morfismo} $\theta:\mathcal{G}\to\mathcal{F}$ (3.5.1) de
haces de anillos (resp. de haces de anillos topológicos).
El \emph{compuesto} de un segundo morfismo
$\Psi'=(\psi',\theta'):(Y,\mathcal{B})\to(Z,\mathcal{C})$ y de $\Psi$, que se denota
por $\Psi''=\Psi'\circ\Psi$, es el morfismo $(\psi'',\theta'')$ en el que
$\psi''=\psi'\circ\psi$, y $\theta''$ es el compuesto de $\theta$ y $\theta'$
(igual a $\psi'_*(\theta)\circ\theta'$, cf. 3.5.2). Para los espacios
anillados, recordemos que entonces se tiene
$\theta''^\sharp=\theta^\sharp\circ\psi^*(\theta'^\sharp)$ (3.5.5). Por tanto, si
$\theta'^\sharp$ y $\theta^\sharp$ son homomorfismos \emph{inyectivos}
(resp. \emph{sobreyectivos}), también lo es $\theta''^\sharp$, pues
$\psi_x\circ\rho_{\psi(x)}$ es un isomorfismo para todo $x\in X$
(3.7.2). De lo anterior se deduce de inmediato que, cuando $\psi$ es una
aplicación continua \emph{inyectiva} y $\theta^\sharp$ es un homomorfismo
\emph{sobreyectivo} de haces de anillos, el morfismo
$(\psi,\theta)$ es un \emph{monomorfismo} (T, 1.1) en la categoría de los
espacios anillados.

Por abuso de lenguaje, se sustituirá a menudo $\psi$ por $\Psi$ en la
notación, por ejemplo escribiendo $\Psi^{-1}(U)$ en lugar de $\psi^{-1}(U)$
para una parte $U$ de $Y$, cuando no haya riesgo de confusión.
\end{env}
